% ============================================================
%  5. fejezet, Összegzés
% ============================================================
\chapter{Összegzés}
\label{ch:osszegzes}

\section{Főbb eredmények}

A KCMamba architektúra differenciálható Kalman-szűrési mechanizmust
integrál a Mamba SSM-be. A három adatkészleten végzett kísérletekből
az alábbi következtetések vonhatók le:

\begin{enumerate}
    \item Zajrobusztusság: $\sigma=0 \to \sigma=5$ zajszint-növelésnél
          a KCMamba átlagos MSE-degrádációja $+1\,528\%$, a Fair Mambáé $+28\,234\%$,
          ami $18{,}5\times$-os különbség. Exchange/$s=16$-nál
          a különbség öt nagyságrend: $+66\%$ vs $+152\,350\%$.

    \item Kalman-gain adaptáció: a modell által az adatokból
          megtanult $K$-érték a zajszint emelkedésével monoton csökken
          ($0.706 \to 0.073$), ami konzisztens az explicit Kalman-egyenletekből
          származtatható várakozásokkal. Ez megerősíti, hogy az architekturális
          prior érdemben irányítja a tanulási folyamatot.

    \item Kevés adat: $s=16$ stride esetén a Kalman-struktúra
          implicit regularizáló hatást fejt ki, ami ellensúlyozza
          az adathiányt.
\end{enumerate}

\section{Korlátok}

A KCMamba nem minden konfigurációban jobb, és a kísérleti keretrendszer
több olyan egyszerűsítést tartalmaz, amelyek a következtetések
általánosíthatóságát érdemben korlátozzák.

\begin{itemize}
    \item Zajtalan, kis stride esetén ($\sigma=0$, $s=16$) a Fair Mamba
          szinte memorizál, és itt az MSE-je sokkal jobb. Ez arra utal,
          hogy az explicit Kalman-struktúra zajtalan környezetben
          felesleges overheadet jelent: a mindig jelen lévő szűrési tag
          enyhe alulilleszkedést okoz, és a struktúra nem képes teljesen
          kikapcsolni a szűrést.

    \item ETTh1/$s=1$-nél az LSTM a legjobb; a gating mechanizmus
          kihasználja a nagy adatmennyiséget. Az LSTM forget/input gate
          párosa hasonló adaptív szűrést valósít meg, de rugalmasabb, mert
          nem köti meg a $K = Q/(Q+R)$ struktúra. Ez nem a KCMamba
          kudarca, hanem annak illusztrációja, hogy adatgazdag rezsimben
          az induktív prior haszna csökken.

    \item Mindössze 5 epoch futott, és minden konfigurációt csak egyetlen
          magon (seed) futtattam. Ez sem a konvergált állapotot, sem a
          statisztikai szignifikanciát nem garantálja: a bemutatott
          eredmények pontbecslések, konfidencia-intervallum nélkül.
          Publikációérett tanulmány legalább három ismétlést és korai
          leállítást igényelne.

    \item Csak $L=H=96$ horizontot vizsgáltam. A 336, 720 és 960 lépéses
          horizontokon a Mamba $O(T)$ komplexitása számottevő előnyt hozhat
          az $O(T^2)$ Transzformerekkel szemben, és az LSTM memória-korlátai
          is itt kezdenek dominálni; a jelen fejezet következtetései
          tehát formálisan csak a rövid horizontra vonatkoznak.

    \item A zajmodell kizárólag izotróp Gauss-zaj
          ($\varepsilon \sim \mathcal{N}(0, \sigma^2 I)$); a valós alkalmazásokban
          a zaj jellemzően heteroszkedasztikus, időben korrelált, esetenként
          nem-Gauss eloszlású, és gyakran strukturális (szenzorkiesés, hiányzó
          érték). Az itt mért robusztusság ezekre a zajformákra nem
          általánosítható automatikusan.

    \item A skalár $Q_t$, $R_t$ becslés feltételezi, hogy a zajkomponensek
          csatornánként függetlenek. Diagonális vagy teljes kovarianciamátrix
          becslése gazdagabb zajmodellt eredményezne, de a paraméterszámot
          érzékelhetően növelné.

    \item Két benchmark (Weather, ECL) a tanítási futamszám skálázódása
          miatt a jelen kísérleti körbe nem került be; a dimenzió-érzékenység
          (7 vs.\ 21 vs.\ 321 csatorna) \emph{a~priori} nincs alátámasztva.

    \item A hiperparaméter-érzékenység szélesebb körű vizsgálata (rejtett
          dimenzió, fejek száma, rétegek száma, kernelméret, tanulási sebesség)
          nem történt meg; a bemutatott konfiguráció egy ésszerű alapbeállítás,
          amely nem biztos, hogy Pareto-optimális.

    \item A valós kereskedelmi oldalon (\ref{ch:implementacio}.\ fejezet)
          a rendszer backtestben van validálva; élő piac, slippage-ingadozás,
          szünetek és részteljesülési sorrend-hatások a jelen dolgozat
          hatókörén kívül esnek.
\end{itemize}

\section{Jövőbeli irányok}

\begin{itemize}
    \item Az $s \in \{1, 4, 8, 16, 32, 64\}$ értékek szisztematikus
          végigpásztázása megmutatná, melyik adatszegénységi küszöbnél
          jelenik meg először a Kalman-struktúra szabályozó hatása.

    \item 336, 720, illetve 960 lépéses előrejelzési horizonton
          a Mamba $O(T \log T)$ komplexitása számottevő előnyt hozhat
          az $O(T^2)$ Transzformerekkel szemben.

    \item A szimulált Gauss-zaj szenzorkiesésekkel, hiányzó értékekkel
          és kilógó megfigyelésekkel való felváltása az architektúrát
          a valós alkalmazási feltételekhez közelítené.

    \item Egy meta-hálózat alapú zajszint-becslés lehetővé tenné,
          hogy $\sigma$ tanított paraméterré váljon, és a modell
          ismeretlen zajú környezetekhez automatikusan alkalmazkodjon.

    \item A skalár $Q_t$, $R_t$ diagonális Cholesky-faktorizációval
          való felváltása gazdagabb bizonytalanság-becslést és
          jobb kalibrációt eredményezhet.
\end{itemize}

\section{Zárszó}

A KCMamba eredményei arra mutatnak rá, hogy a klasszikus Kalman-szűrési
elvek és a modern szekvencia-modellek nem egymást kizáró megközelítések,
hanem produktívan összekapcsolhatók.
Az explicit zajmodell olyan induktív torzítást visz az architektúrába,
amely előnyt jelent a korlátozott adatmennyiségű és/vagy magas
zajszintű körülmények között.
A~\ref{ch:implementacio}. fejezetben ismertetett rendszer azt is
igazolja, hogy a modell valós kereskedési környezetbe integrálható,
nem csak kísérleti körülmények között használható.
